\documentclass[11pt]{article}
\usepackage[T1]{fontenc}
\usepackage[utf8]{inputenc}
\usepackage[english]{babel}
\usepackage{lmodern}
\usepackage{amsmath,amssymb,mathtools,array,booktabs}
\usepackage{geometry}
\usepackage{microtype}
\geometry{a4paper,margin=1.45cm}
\setlength{\parindent}{1.25em}
\setlength{\parskip}{0pt}
\allowdisplaybreaks
\setlength{\emergencystretch}{30em}
\sloppy
\hbadness=10000
\vbadness=10000
\hfuzz=10pt
\vfuzz=10pt
\raggedbottom
\providecommand{\mG}{\mathfrak{G}}
\providecommand{\mA}{\mathfrak{A}}
\providecommand{\mB}{\mathfrak{B}}
\providecommand{\mC}{\mathfrak{C}}
\providecommand{\mD}{\mathfrak{D}}
\providecommand{\mE}{\mathfrak{E}}
\providecommand{\mH}{\mathfrak{H}}
\providecommand{\mK}{\mathfrak{K}}
\providecommand{\mM}{\mathfrak{M}}
\providecommand{\mN}{\mathfrak{N}}
\providecommand{\mR}{\mathfrak{R}}
\providecommand{\mS}{\mathfrak{S}}
\providecommand{\mT}{\mathfrak{T}}
\providecommand{\mU}{\mathfrak{U}}
\providecommand{\mO}{\mathfrak{O}}
\providecommand{\mo}{\mathfrak{o}}
\providecommand{\OO}{\Omega}
\providecommand{\Th}{\Theta}
\newcommand{\iso}{\simeq}
\newcommand{\normlhd}{\triangleleft}
\newcommand{\mat}[1]{\begin{pmatrix}#1\end{pmatrix}}
\begin{document}

\providecommand{\mG}{\mathfrak{G}}
\providecommand{\mA}{\mathfrak{A}}
\providecommand{\mB}{\mathfrak{B}}
\providecommand{\mC}{\mathfrak{C}}
\providecommand{\mD}{\mathfrak{D}}
\providecommand{\mE}{\mathfrak{E}}
\providecommand{\mH}{\mathfrak{H}}
\providecommand{\mK}{\mathfrak{K}}
\providecommand{\mM}{\mathfrak{M}}
\providecommand{\mN}{\mathfrak{N}}
\providecommand{\mR}{\mathfrak{R}}
\providecommand{\mS}{\mathfrak{S}}
\providecommand{\mT}{\mathfrak{T}}
\providecommand{\mU}{\mathfrak{U}}
\providecommand{\mO}{\mathfrak{O}}
\providecommand{\mo}{\mathfrak{o}}
\providecommand{\OO}{\Omega}
\providecommand{\Th}{\Theta}
\providecommand{\iso}{\simeq}
\providecommand{\normlhd}{\triangleleft}

\setcounter{footnote}{0}

\section*{34. Hypercomplex Quantities and Representation Theory}
\emph{Math. Zs. 30 (1929), pp. 641--692.}

\subsection*{Introduction}

The most important general theorems on hypercomplex systems go back to Molien (Math. Annalen vol. 41 and the Dorpat reports of 1897). Essentially independently of this, Frobenius soon afterward developed the theory of hypercomplex systems and of their representations -- in particular the representation theory of finite groups -- as a unified theory. The basis is the concept, due to Dedekind, of the group determinant, more generally the determinant of an arbitrary hypercomplex system. Frobenius shows that to the various irreducible factors of this determinant (where the coefficient domain is always the field of all complex numbers) there correspond the various irreducible representation classes, and that in this way all irreducible representation classes are exhausted. Such a representation class is completely characterized by its ``character system''; in the case of finite groups this character system arises by decomposing the determinant of the commutative hypercomplex system derived from the classes of conjugate elements of the group. This determinant decomposes into linear factors, with the characters as coefficients -- a direct generalization of the result first obtained by Dedekind, namely that the group determinant of a finite Abelian group decomposes into linear factors whose coefficients are the different characters of the Abelian group (correspondence with Frobenius).\footnote{What follows is a free elaboration, prepared by B. L. van der Waerden, of my lectures of the winter semester 1927/28. We prepared the text for print jointly. I am also indebted to B. L. van der Waerden for a number of critical remarks.}

These conceptually simple and transparent results, however, are obtained by Frobenius through laborious calculation. The later development aimed at a simplified derivation of these results, and at the same time at their extension when an arbitrary field is taken as coefficient domain. This later development proceeded in completely separate ways for hypercomplex systems and for representation theory. Hypercomplex systems received an arithmetical treatment through Wedderburn: every hypercomplex system is uniquely a direct sum of two-sided directly indecomposable rings; for systems without radical these indecomposable rings become isomorphic to the system of all $n$-rowed matrices with elements from an assigned, not necessarily commutative, division ring -- determined essentially uniquely.\footnote{Cf. the references in Dickson, \emph{Algebras and their Arithmetics} (German edition: \emph{Algebren und ihre Zahlentheorie}, Zürich 1927).} Representation theory received an elementary foundation through Burnside and I. Schur, who proceeded directly from a given representation, independently of the hypercomplex system, and operated with matrix theorems. In particular Schur treated the question of the number fields of smallest degree in which a representation irreducible over a given field decomposes absolutely. These absolutely irreducible constituents belong to finitely many conjugate representation classes. The degree of the number fields sought is equal to the product of the number of these classes by the index (the number of constituents in each of the conjugate classes). The chief auxiliary tool of Schur's theory is the system of all matrices commuting with an irreducible representation.\footnote{Cf. the references in Chapter 10 of Speiser's group theory.}

In the following purely arithmetical foundation, hypercomplex quantities and representation theory again appear as one unified whole -- as a special case of a general theory of non-commutative rings satisfying only certain finiteness conditions. More precisely, this is a theory of module and ideal classes with respect to these rings, with the main result that the irreducible module classes are already exhausted by the corresponding ideal classes; in particular, for rings without radical all module classes decompose into irreducible ones (become completely reducible). This is the arithmetical equivalent of Frobenius's result that the irreducible factors of the regular group (respectively system) determinant already exhaust the totality of irreducible representation classes, and that complete reducibility of representations holds for systems without radical. Namely, the consideration of the system determinant and its decomposition can be interpreted as a transition to the norm, whereas here the direct-sum decomposition of the ideals and their composition series are treated directly. Representation theory, by means of the concept of the representation module, becomes a theory of module classes.

The introduction of module and ideal classes has a purely group-theoretic basis. Modules and ideals are regarded as Abelian groups with respect to addition, restricted by the condition that they admit certain multiplications by ring elements: they form ``groups with operators'' (§1). For such groups, in place of the ordinary isomorphism there appears the ``operator isomorphism'' (§2); groups that are operator-isomorphic to one another (within a fixed domain) are collected into a class -- a class concept that, for example in the case of the ideals of a number field, coincides with the usual one. The theory of groups with operators goes back to W. Krull and O. Schmidt (cf. note 6); it is developed systematically in Chapter I. The theorems that remain valid here as well -- on composition series, direct product (respectively sum), and completely reducible groups; uniqueness theorems in the sense of the class division just mentioned -- form the basis of everything that follows.

They first yield the general uniqueness theorems for the irreducible diagonal constituents of arbitrary representations (Chapter III, §16), on the basis of the one-to-one correspondence of representation classes with the classes of representation modules, hence with classes of groups with operators. The further question of the totality of representations requires a more precise study of the structure of the rings to be represented, thus, in the special case, of hypercomplex systems. In Chapter II the results of Wedderburn are obtained anew and carried further, on the basis of the group-theoretic conception according to which one-sided ideals stand in the foreground. It turns out that the ``multiple-chain theorem'' for right ideals, or the identical ``minimal condition'' (in every set of right ideals there is at least one minimal member within the set), suffices as a finiteness condition.\footnote{Wedderburn's methods of proof can be transferred if the ``double chain theorem'' is assumed. Cf. E. Artin, Zur Theorie der hyperkomplexen Zahlen, Hamb. Abh. 1927. Cf. also A. Suschkewitsch, Über die endlichen Gruppen ohne das Gesetz der eindeutigen Umkehrbarkeit, Math. Annalen 99 (1928), pp. 30--50, where parallel structure theorems are developed for (finite) domains with only one associative, non-invertible operation. The splitting of the ``kernel'' in Suschkewitsch into right and left groups corresponds to the representation of a two-sided simple ring with identity element as a sum of right and left ideals.} One obtains the identity of the rings without radical satisfying the minimal condition with the (right) completely reducible rings with identity element. In such rings all simple right ideals of a two-sided indecomposable ring belong to the same class and therefore possess -- as groups with operators -- the same automorphism ring, which, because of the simplicity of the ideals, becomes a division ring. This automorphism division ring of the class is what mediates the matrix representation found by Wedderburn.

This matrix representation in the automorphism division ring also solves the representation problem in the completely reducible case, as soon as representation with respect to the automorphism division ring or to its subfields is required -- in particular with respect to the coefficient domain of a hypercomplex system. There are -- this is a direct consequence of the theorem reducing module classes to ideal classes -- no further irreducible representations beyond the Wedderburn representation and those that arise from it when the elements of the automorphism division ring are replaced in the natural way by matrices over a subfield.

Thus in the completely reducible case there are as many different irreducible representation classes as ideal classes; their number agrees with the number of different two-sided indecomposable, hence simple, rings. This number is finally again identical with the number of indecomposable components of the center, which become commutative fields. In the case of hypercomplex systems with algebraically closed coefficient domain, these latter fields are completely determined by the different ring homomorphisms (irreducible representations) of the center; up to a numerical factor these are precisely the characters.

The results obtained also give, at the same time, a survey of the irreducible representations of systems with radical, insofar as these can be reduced to the representations of the residue class ring modulo the radical, which becomes a system without radical.

As will be shown later, for hypercomplex systems without radical the theory of ``splitting fields'' also follows from the automorphism division ring: that is, the theory of commutative extensions of the coefficient domain in which the representations decompose into absolutely irreducible representations -- equivalently, in which a direct-sum decomposition into absolutely simple right ideals occurs. The splitting fields are identical with those of the automorphism division ring, and all splitting fields of smallest degree are isomorphic to the maximal commutative subfields of the automorphism division ring. The connection with Schur's investigations mentioned above is established by the fact that the transposes of the matrices commuting with the representation supply precisely a representation of the automorphism division ring.\footnote{Cf. R. Brauer--E. Noether, Über minimale Zerfällungskörper irreduzibler Darstellungen, Sitz.-Ber. Preuß. Akad. Wiss. 1927, p. 221.} These theorems are to be developed systematically in the framework of a Galois theory of non-commutative division rings.

The underlying concepts -- operator homomorphism and automorphism ring -- also yield the structure of general rings with radical that satisfy only the minimal condition; this will be carried out from another side.

\section*{Chapter I. Group-Theoretic Foundations}

\subsection*{§ 1. Groups with operators}

Let $\mG$ be a group (finite or infinite), with elements $a,b,\ldots$. By an operator domain $\Omega$ for the group $\mG$ is meant a set of new symbols $H,\Theta,\ldots$, such that to every $a$ in $\mG$ and every $\Theta$ in $\Omega$ there is assigned a uniquely defined $\Theta a$ in $\mG$, and such that the distributive law holds:
\[
        \Theta(ab)=\Theta a\cdot\Theta b.
\]
Accordingly, every operator defines a homomorphism of the group into itself, where the group is mapped either onto itself or onto a subgroup. The identity element goes into the identity element, and inverses go into inverses.

An operator domain is called \emph{absolute} if different operators also define different homomorphisms. From an arbitrary operator domain one obtains an absolute one by equating all those elements which generate the same homomorphism. An absolute operator domain is a one-to-one image of a subset of the set of all homomorphisms of the group into itself.

An \emph{admissible subgroup} $\mH$ of $\mG$ -- with respect to a fixed operator domain $\Omega$ -- is a subgroup which admits the operator domain, that is, for which $\Theta a$ lies in $\mH$ for every $a$ in $\mH$ and $\Theta$ in $\Omega$.

\emph{Examples.} 1. Let the operators be the inner automorphisms: $\Theta a=c^{-1}ac$. The admissible subgroups are the normal divisors.

2. Let the operators be all automorphisms. The admissible subgroups are the ``characteristic subgroups'', which go into themselves under all automorphisms.\footnote{The concepts explained here come from W. Krull, Über verallgemeinerte endliche Abelsche Gruppen, Math. Zeitschr. 23 (1925), pp. 161--196, and O. Schmidt, Über unendliche Gruppen mit endlicher Kette, Math. Zeitschr. 29 (1928), pp. 34--41.}

3. Let $\mG$ be a ring, i.e. an Abelian group with respect to addition, in which a multiplication is also defined, with the properties
\[
        r(a+b)=ra+rb,\qquad (a+b)r=ar+br,\qquad ab\cdot c=a\cdot bc.
\]
Every element $r$ simultaneously defines two operators: the operators $rx$ and $xr$. The admitted subgroups are the ``ideals'' $\mathfrak a$, namely: left-sided ideals, which admit the operations $rx$: $r\mathfrak a\subseteq\mathfrak a$; right-sided ideals, which admit the operations $xr$: $\mathfrak a r\subseteq\mathfrak a$; and two-sided ideals, which admit both operations. All ideals become trivial $(=\{0\}\text{ or }=\mG)$ when the ring $\mG$ is a division ring, i.e. when $\mG-\{0\}$ is a group with respect to multiplication.\footnote{Thus neither for rings nor for division rings is the commutative law of multiplication assumed.}

4. \emph{Modules with respect to a ring $\mo$.} Let $\mo$ be a ring, and $\mM$ an Abelian group written additively. Suppose a multiplication $r\cdot a$ of elements of $\mo$ with elements of $\mM$ is given; the product is to lie again in $\mM$. One requires:
\[
\left.
\begin{aligned}
 r(a+b)&=ra+rb,\\
 rs\cdot a&=r\cdot sa,\\
 (r+s)a&=ra+sa,
\end{aligned}
\right\}\qquad r,s\in\mo,\quad a,b\in\mM.
\]
Then $\mM$ is called an $\mo$-module; more precisely, since the multipliers from $\mo$ are written on the left, a left $\mo$-module. The ring $\mo$ is at the same time an operator domain. If it is absolute, i.e. if different ring elements also yield different operations, one speaks of an \emph{absolute multiplier domain} for the module $\mM$.

As above, one can pass from an arbitrary multiplier domain to the absolute one by equating the elements which generate the same operation. The absolute multiplier domain is again a ring. The admissible subgroups of a module $\mM$ are called \emph{submodules}. If, in particular, $\mo=\mM$, then one obtains the left ideals again.\footnote{In the same way one requires for right modules:
\[
(a+b)r=ar+br,\qquad a\cdot rs=ar\cdot s,\qquad a(r+s)=ar+as.
\]}

5. \emph{Bimodules.} If $\mM$ is simultaneously a left $\mo$-module and a right $\mo'$-module, and if moreover for $a$ in $\mo$, $\alpha$ in $\mM$, and $a'$ in $\mo'$ one always has
\[
        a\cdot\alpha a'=a\alpha\cdot a',
\]
then $\mM$ is called a bimodule.

6. \emph{The automorphism ring of an Abelian group.}\footnote{Cf. A. Châtelet, \emph{Les Groupes Abéliens finis}, Paris, Gauthier-Villars, 1925, p. 99.} For the endomorphisms of an Abelian group (written additively) one can define an addition and a multiplication by the formulas
\[
        (H+\Theta)a=Ha+\Theta a,\qquad (H\Theta)a=H(\Theta a).
\]
The operations so defined plainly again represent homomorphisms, and therefore again belong to the system. The system forms a ring, and the Abelian group can, according to 4, be regarded as a module with respect to this ring.

In what follows, by ``subgroups'' we shall always mean admissible subgroups. The intersection $\mA\cap\mB$ of two admissible subgroups is again an admissible subgroup. So is the product $\mA\mB$, provided at least one of the two subgroups $\mA,\mB$ is a normal divisor.

\subsection*{§ 2. The isomorphism theorems}

A mapping of a group $\mG$ onto a group $\overline{\mG}$ -- where $\mG$ and $\overline{\mG}$ are to possess the same operator domain -- is called an operator homomorphism if, first, it is a homomorphism in the ordinary sense ($ab\mapsto \bar a\bar b$), and if, second, whenever $a$ maps to $\bar a$, the element $\Theta a$ maps to $\Theta\bar a$.\footnote{One may extend the concept of operator homomorphism by allowing the groups $\mG,\overline{\mG}$ to have separate operator domains, the homomorphism then mapping not only $\mG$ onto $\overline{\mG}$ but also the operator domains onto one another, in such a way that, if $a$ maps to $\bar a$ and $\Theta$ maps to $\overline\Theta$, then $\overline\Theta\bar a$ is the image of $\Theta a$. In this general form the concept of operator homomorphism also includes that of ring homomorphism; cf. §8.} Notation: $\mG\sim\overline{\mG}$. If the correspondence is one-to-one, it is called an operator isomorphism. Notation: $\mG\simeq\overline{\mG}$. As for ordinary groups, one proves the homomorphism theorem:

\emph{If $\overline{\mG}$ is an operator-homomorphic image of $\mG$, then $\overline{\mG}$ is operator-isomorphic to a factor group $\mG/\mN$, where $\mN$ is an admissible normal divisor, consisting of all elements of $\mG$ that correspond to the identity in $\overline{\mG}$. Conversely, every admissible normal divisor $\mN$ defines a factor group $\mG/\mN$ which admits the operator domain of $\mG$ and is an operator-homomorphic image of $\mG$.}

A group is called \emph{simple} if it has no normal divisors other than itself and the identity group. Every homomorphism of a simple group is either an isomorphism, or else assigns the identity element to every element.

\emph{First isomorphism theorem.} \emph{Let $\overline{\mG}$ be a homomorphic image of $\mG$, let $\overline{\mA}$ be a normal divisor of $\overline{\mG}$, and let $\mA$ be the totality of those elements of $\mG$ to which elements of $\overline{\mA}$ are assigned. Then $\mA$ is again a normal divisor, and}
\[
        \overline{\mG}/\overline{\mA}\simeq \mG/\mA. \tag{1}
\]
\emph{Proof.} $\mG\sim\overline{\mG}$, $\overline{\mG}\sim\overline{\mG}/\overline{\mA}$, hence $\mG\sim\overline{\mG}/\overline{\mA}$. Therefore $\overline{\mG}/\overline{\mA}$ is isomorphic to a factor group in $\mG$; the corresponding normal divisor is the totality of elements that correspond to the identity in $\overline{\mG}/\overline{\mA}$; this is precisely $\mA$.

\emph{Addendum.} If, according to the homomorphism theorem, one sets $\overline{\mG}=\mG/\mN$, then $\mA$ certainly contains $\mN$. From $\mA$ one can recover $\overline{\mA}$: $\overline{\mA}=\mA/\mN$. The assignment $\mA\leftrightarrow\overline{\mA}$ is one-to-one. Formula (1) can also be written
\[
        (\mG/\mN)/(\mA/\mN)\simeq \mG/\mA.
\]

\emph{Second isomorphism theorem.} \emph{Let $\mA$ be a subgroup and $\mB$ a normal divisor of $\mG$. Then $\mA\cap\mB$ is a normal divisor in $\mA$, and}
\[
        \mA\mB/\mB\simeq \mA/(\mA\cap\mB).
\]
\emph{Proof.} Under the homomorphism $\mG\sim\mG/\mB$, in particular the elements $a$ of $\mA$ are assigned certain residue classes $a\mB$, which together make up the group $\mA\mB/\mB$. Thus $\mA\sim\mA\mB/\mB$, from which the assertion follows by the homomorphism theorem.

An operator-homomorphic mapping of a group $\mG$ onto itself or onto a subgroup is called an endomorphism of $\mG$. If $\mG$ is in particular an Abelian group (with operators), written additively, and if sum and product of homomorphisms are defined as above (§1, Example 6), then the operator homomorphisms of the group into itself form a ring, the automorphism ring.

\emph{The automorphism ring of a simple Abelian group (with operators) is a division ring.}

\emph{Proof.} Every homomorphism maps the group either onto itself or onto the zero group (since there are no other admissible subgroups). The homomorphisms that do not map everything to zero are, by what was remarked above about simple groups, isomorphisms, and the isomorphic mappings of a group onto itself form a group. Thus after omitting the zero operator the ring becomes a group; hence the ring itself is a division ring.

If, in particular, $\mG$ is a left $\mo$-module, and if the operator homomorphisms are written as right operators, then $\mG$ becomes a \emph{bimodule with respect to $\mo$ on the left and the automorphism ring on the right}. For the fact that the new operators $\Gamma$ were chosen as homomorphisms is expressed by the formulas
\[
        (a+b)\Gamma=a\Gamma+b\Gamma,
        \qquad ra\cdot\Gamma=r\cdot a\Gamma,
\]
which (together with the remaining formulas already interpreted earlier) characterize the bimodule.

Conversely, if a bimodule is given, then these same formulas express that every element of the right multiplier domain induces an operator homomorphism with respect to the left operators, and conversely.

\subsection*{§ 3. Composition series}

If in $\mG$ there is a finite sequence of admissible subgroups
\[
        \mG>\mA_1>\mA_2>\cdots>\mA_r=\mE \tag{2}
\]
($\mE$ is the group consisting of the identity element $e$ alone), such that every $\mA_i$ is a normal divisor in the preceding member, and it is impossible to insert between two successive members of the sequence another subgroup with the same property, then the sequence is called a composition series.

The factor groups $\mG/\mA_1,\ \mA_1/\mA_2,\ldots,\ \mA_{r-1}/\mE$ are called composition factors. They are simple groups, i.e. they possess no admissible normal divisors other than themselves and the identity group.

If among the operators one includes, in particular, all inner automorphisms, then the sequence (2) consists entirely of normal divisors, and one calls it a principal series. If one includes all automorphisms, it is called a characteristic series. In the later examples of §1 this would mean composition series of ideals in $\mo$, respectively composition series of modules or bimodules.

\emph{Jordan--Hölder theorem.} \emph{If for a group $\mG$ two different composition series exist,}
\[
        \mG>\mA_1>\mA_2>\cdots>\mA_r=\mE,
        \qquad
        \mG>\mB_1>\mB_2>\cdots>\mB_s=\mE,
\]
\emph{then they have the same length, $r=s$, and the composition factors}
\[
        \mG/\mA_1,\ \mA_1/\mA_2,\ldots,\mA_{r-1}/\mE
\]
\emph{are, in some order, isomorphic to}
\[
        \mG/\mB_1,\ \mB_1/\mB_2,\ldots,\mB_{s-1}/\mE.
\]
For the proof, see for example E. Noether, Abstrakter Aufbau usw., Math. Annalen 96 (1926), §10, p. 57. At the same time it is proved there that:

\emph{If in $\mG$ there is a composition series, then a composition series can be drawn through every normal divisor $\mH$ of $\mG$.}

The existence of a composition series is clear for finite groups, more generally for those groups in which a maximal condition and a minimal condition hold:

The \emph{maximal condition} says that in every set of subgroups there is a maximal subgroup, i.e. one that is no longer contained in another subgroup of the set. Equivalently, every ascending chain of subgroups $\mA_1\subset\mA_2\subset\mA_3\cdots$ breaks off after finitely many members.

The \emph{minimal condition} says that in every set of subgroups there is a minimal subgroup, one that contains no other subgroup of the set; equivalently, every descending chain $\mA_1\supset\mA_2\supset\mA_3\cdots$ breaks off after finitely many members.

\emph{If only normal divisors are admitted as subgroups} (for example for Abelian groups), \emph{then conversely the maximal and minimal conditions follow from the existence of the composition series.}

\emph{Proof.} Every normal divisor has a composition series, whose length is called the length of the normal divisor. If $\mA\subset\mB$, then the length of $\mA$ is smaller than that of $\mB$, since a composition series for $\mB$ can be drawn through $\mA$. Thus every subgroup of shortest length is at the same time minimal, and every subgroup of greatest length is at the same time maximal.

\subsection*{§ 4. Direct products and intersections}

A group $\mG$ is called the direct product of two factors, $\mG=\mA\times\mB$, if
\[
\begin{array}{ll}
1.& \mA,\mB\text{ are normal divisors in }\mG,\\
2.& \mA\mB=\mG,\\
3.& \mA\cap\mB=\mE.
\end{array}
\]
\emph{The definition is plainly equivalent to the following:} Every element $g$ of $\mG$ has a unique representation as $g=ab$, with $a$ in $\mA$ and $b$ in $\mB$, and the elements of $\mA$ commute with those of $\mB$: $ab=ba$.

A group $\mG$ is called the direct product of $n$ factors, $\mG=\mA_1\times\cdots\times\mA_n$, if, putting $\mB_i=\mA_1\cdots\mA_{i-1}\mA_{i+1}\cdots\mA_n$, one has $\mG=\mA_i\times\mB_i$ directly for every $i$.

\emph{The definition is again equivalent to the following:} Every $g$ of $\mG$ is uniquely representable as
\[
        g=a_1a_2\cdots a_n,\qquad a_i\in\mA_i,
\]
and the elements of $\mA_i$ commute with those of $\mA_k$.

The following facts are often used:
\begin{enumerate}
\item If $\mA_1\times\cdots\times\mA_n=\mR$ and $\mR\times\mH=\mG$, then $\mA_1\times\cdots\times\mA_n\times\mH=\mG$.
\item If $\mG=\mA_1\times\cdots\times\mA_n=\mC_1\times\cdots\times\mC_n$ and $\mC_i\subseteq\mA_i$, then $\mC_i=\mA_i$. (This follows from the representation of the elements of $\mA_i$ by means of the $\mC_j$.)
\item If $\mG=\mA\times\mB$ and $\mR\supseteq\mA$, then $\mR=\mA\times(\mR\cap\mB)$. (This follows from the representation of the elements of $\mR$ in the form $ab$, where the second factor belongs both to $\mR$ and to $\mB$.)
\item If $\mG=\mA\times\mB$, then $\mA\sim\mG/\mB$ and $\mB\sim\mG/\mA$ (second isomorphism theorem). It follows further that:
\item If $\mG=\mA\times\mB$, and if $\mA$ and $\mB$ possess composition series of lengths $m$ and $n$, then $\mG$ possesses a composition series of length $m+n$.
\item If $\mA\sim\overline{\mA}$ and $\mB\sim\overline{\mB}$, then $\mA\times\mB\sim\overline{\mA}\times\overline{\mB}$.
\end{enumerate}
A group is called directly indecomposable if it cannot be represented as a direct product of factors $\ne\mE$.

It is clear that \emph{every group satisfying the minimal condition is the direct product of finitely many directly indecomposable groups}.\footnote{As W. Krull in the commutative case and O. Schmidt in general have shown (cf. note 6), under the hypothesis of the maximal and minimal conditions the representation is unique up to isomorphism. Since one can add all inner automorphisms to the operator domain without changing the concept of direct indecomposability, it is enough to require the finiteness conditions for normal divisors, or the existence of a principal series.}

If the groups in question are written additively, then the notions product and direct product pass into sum and direct sum. Notation: for a sum of groups, $(\mA_1,\mA_2,\ldots)$; for a direct sum, $\mA_1+\mA_2+\cdots$.

The concept of ``direct intersection'' will not be used later, but the following theorems concerning the connection between direct intersection and direct product show how the ideal theory of the next chapters can also be formulated with intersections instead of sums, thereby restoring the connection with the familiar commutative theories.

A group $\mD$ is called the direct intersection of two groups $\mR$ and $\mS$ with respect to $\mG$ if
\[
\begin{array}{ll}
1.& \mR,\mS\text{ are normal divisors in }\mG,\\
2.& \mR\cap\mS=\mD,\\
3.& \mR\mS=\mG.
\end{array}
\]
Likewise $\mD$ is called the direct intersection of $n$ groups $\mS_1,\ldots,\mS_n$ if, putting $\mR_i=\mS_1\cap\cdots\cap\mS_{i-1}\cap\mS_{i+1}\cap\cdots\cap\mS_n$, the intersection $\mD=\mR_i\cap\mS_i$ is direct for every $i$.

From the definition it follows that $\mD$ must be a normal divisor in $\mG$. If $\mD=\mS_1\cap\cdots\cap\mS_n$ is direct and, under the homomorphism $\mG\sim\mG/\mD$, the groups $\mG,\mS,\mD$ pass to $\overline{\mG},\overline{\mS},\mE$, then $\mE=\overline{\mS}_1\cap\cdots\cap\overline{\mS}_n$ becomes direct. Conversely, from every such representation $\mE=\overline{\mS}_1\cap\cdots\cap\overline{\mS}_n$ one returns to a direct representation $\mD=\mS_1\cap\cdots\cap\mS_n$. Through this assignment theorems on direct intersections for arbitrary $\mD$ are reduced to the special case $\mD=\mE$.

In the case $\mD=\mE$ and two factors, the conditions for direct product and direct intersection coincide. For $n$ factors there is the following one-to-one relation between direct product and intersection:

\emph{I. If $\mG=\mA_1\times\cdots\times\mA_n=\mA_i\times\mB_i$, then $\mE=\mB_1\cap\cdots\cap\mB_n=\mB_i\cap\mC_i$ directly, and $\mC_i=\mA_i$.}

\emph{II. If $\mE=\mS_1\cap\cdots\cap\mS_n=\mR_i\cap\mS_i$ directly, then $\mG=\mR_1\times\cdots\times\mR_n=\mR_i\times\mT_i$, and $\mT_i=\mS_i$.}

\emph{Proof of I.} Let $\mB=\mB_1\cap\cdots\cap\mB_n=\mB_i\cap\mC_i$; then $\mC_i\supseteq\mA_i$. We show $\mC_i\subseteq\mA_i$. If, for instance, $c$ lies in $\mC_1$, then $c$ lies in $\mB_2,\ldots,\mB_n$, so that the component representation with respect to the $\mA$ gives
\[
        c=a_1a_2\cdots a_n=a_1ea_3\cdots a_n=a_1a_2e\cdots a_n=\cdots=a_1\cdots a_{n-1}e.
\]
Since the product of all $\mA_i$ is direct, the representations agree; in every position except the first there is once an $e$, hence $c=a_1$, or $\mC_1\subseteq\mA_1$, i.e. $\mC_1=\mA_1$, and correspondingly $\mC_i=\mA_i$. It follows that $\mB=\mB_i\cap\mC_i=\mB_i\cap\mA_i=\mE$; moreover $\mB_i\mC_i=\mB_i\mA_i=\mG$, hence the intersection is direct.

\emph{Proof of II.} Let $\mH=\mR_1\cdots\mR_n=\mR_i\mT_i$; then $\mT_i\subseteq\mS_i$. We show $\mS_i\subseteq\mT_i$. Let, for instance, $s$ lie in $\mS_1$. By means of $\mG=\mS_i\times\mR_i$ one obtains
\[
        s=s_1e=s_2r_2=\cdots=s_nr_n.
\]
Forming $t_1=er_2\cdots r_n=t_ir_i$, and taking account of the elementwise commutativity of $\mR_i$ and $\mS_i$, one gets
\[
        st_1^{-1}=s_it_i^{-1},
\]
which is therefore an element of all the $\mS_i$, and hence is equal to $e$. Thus $\mS_1\subseteq\mT_1$, or $\mT_1=\mS_1$, and correspondingly $\mT_i=\mS_i$. It follows that $\mH=\mR_i\mT_i=\mR_i\mS_i=\mG$ and $\mR_i\cap\mT_i=\mR_i\cap\mS_i=\mE$; hence $\mG$ is the direct product of the $\mR_i$.

\end{document}
